\section{Discussion}

\subsection{Empirical low-dimensionality of ocean carbon}
The central finding of this study is that the empirical dependence of ocean $\tco$ on salinity, temperature, and $\aou$ is substantially compressible into a low-dimensional nonlinear representation. Across all four ocean basins, a five-parameter rank-1 quadratic model captures the vast majority of predictive structure, outperforming linear regression models while achieving accuracy comparable to a full 10-parameter second-order polynomial.

\subsection{Significance of the rank-1 versus rank-2 comparison}
The comparison between rank-1 and rank-2 quadratic models provides one of the strongest structural insights of our analysis. Increasing the quadratic rank from one to two introduces four additional parameters but produces virtually no predictive benefit ($\Delta\mathrm{RMSE} = +1.98, +0.25, -0.29, -0.04\,\mu\mathrm{mol}\,\mathrm{kg}^{-1}$). In the Atlantic and Indian Oceans, rank-2 models actually degrade out-of-sample precision due to parameter overfitting. This result demonstrates that the quadratic curvature governing empirical $\tco$ distributions is effectively one-dimensional: a single linear projection $z = \alpha S + \beta T + \gamma \aou_{100} + \delta$ accounts for virtually all non-additive quadratic variability.

\subsection{Role of rank-1 as a complexity-efficient representation}
While rank-1 provides a compact parameterization, higher-complexity symbolic regression expressions achieve lower absolute RMSE ($3$--$12\%$ reduction) in all basins. However, these gains require expression complexities of $C=15$--$25$ (compared to $C=8$ for rank-1). Thus, rank-1 should not be claimed as the single "most accurate" model or an invariant "Pareto knee." Instead, rank-1 represents a highly efficient low-complexity approximation that trades a minor fraction of predictive accuracy for exceptional parameter compactness and structural transparency.

\subsection{Operator-space dependence in symbolic recovery}
Our multi-search design ($48$ PySR runs across $4$ operator sets and $3$ seeds) demonstrates that symbolic regression outputs are inherently sensitive to the predefined operator space. Although rank-1 candidates recurred in $24$ of $48$ search cells and appeared in all four basins, recovery in the Atlantic Ocean occurred exclusively when explicit squaring operators were enabled (Search C). In basic arithmetic or log-exponential spaces, genetic algorithms failed to isolate the rank-1 form in the Atlantic. This highlights that symbolic candidate discovery is an interaction between observational data structure and search space parameterization, rather than an automated proof of equation uniqueness.

\subsection{Domain dependence and generalization boundaries}
Evaluating models across temporal validation, spatial-block ($5^\circ \times 10^\circ$) cross-validation, cruise-block cross-validation, and locked post-2018 external holdouts reveals that predictive performance is strongly domain-dependent. Model rankings shift between temporal and spatial validation schemes, emphasizing that no single model family is universally superior across all validation protocols.

Crucially, identifying explicit failure regimes strengthens the scientific utility of compact parameterizations:
\begin{enumerate}[label=(\roman*)]
\item \textbf{Surface Waters ($0$--$100$\,m)}: High predictive error across all model families reflects unmodeled seasonal air--sea $\mathrm{CO}_2$ gas exchange, thermal warming/cooling, and rapid biological uptake.
\item \textbf{Antarctic Intermediate Water (AAIW)}: Severe systematic bias ($+27.14\,\mu\mathrm{mol}\,\mathrm{kg}^{-1}$) in Atlantic AAIW indicates that $S, T, \aou$ alone lack sufficient state information to constrain intermediate water ventilation histories.
\item \textbf{Southern Ocean Abyss ($>3000$\,m)}: Global model failure ($R^2 = -0.174$) contrasted with excellent local refit accuracy ($R^2 = 0.786$) demonstrates that basin-scale parameterizations cannot be extrapolated across distinct deep circulation regimes without local coefficient recalibration.
\end{enumerate}

\subsection{Physical interpretation versus empirical fitting}
Mathematically, the rank-1 Hessian $H = 2 \mathbf{v} \mathbf{v}^T$ establishes a one-dimensional curvature structure in state space. Oceanographically, the monotonic partial derivatives ($\partial\tco/\partial S > 0, \partial\tco/\partial T < 0, \partial\tco/\partial\aou > 0$) align qualitatively with carbonate buffering dynamics and biological remineralization stoichiometry \citep{Sarmiento2006}. However, we explicitly emphasize that empirical curvature alignment does not constitute proof of a fundamental mixing axis, a first-principles derivation of aqueous carbonate equilibrium, or a universal physical law. The rank-1 quadratic is an empirical projection that efficiently captures observational correlations.

\subsection{Limitations}
Several limitations must be acknowledged:
\begin{enumerate}[label=(\roman*)]
\item Symbolic regression searches were restricted to static operator sets and fixed complexity caps ($C_{\max}=30$).
\item Water-mass definitions relied on simplified depth and property thresholds rather than full neutral-density or multi-variable water-mass endmember analysis.
\item Spatial and cruise blocking evaluated existing GLODAP sampling geometries but cannot replace fully independent autonomous sensor networks (e.g., Biogeochemical-Argo floats).
\end{enumerate}

\subsection{Implications for compact ocean carbon parameterizations}
Despite these limitations, the compact rank-1 quadratic parameterization offers practical utility for oceanographic applications. With only five fitted coefficients per basin, it provides an interpretable, computationally lightweight baseline for gap-filling hydrographic bottle data, performing regional carbon diagnostics, and validating complex ocean biogeochemical models without the risk of unconstrained polynomial divergence.
